\section{The finite-alphabet tower before pruning}
\label{tow:section}

This section specifies the vertices, slots, and labels.  The choice of the
permutation tables needed for expansion is made in
Proposition~\ref{exp:realized-gap}.  Section~\ref{prune:section} then makes
the deletions.  Throughout, a bit stream means an element of
\(\mathbb F _2^{\mathbb N_0}\) with a specified finite set of active
coordinates; every inactive coordinate is zero.  A prefix of length
\(a\) comprises coordinates \(0,\ldots,a-1\).  In particular, references
to a prefix longer than the active stream refer to its zero padding.

The construction separates two requirements that are awkward to impose
on one graph at once. An auxiliary graph \(Y^0\) will have expanding
bipartite blocks between neighboring layers. A graph \(Y^\#\) on the same
vertices will have a single labeling rule from which long path words
recover their vertices. Deleting a uniformly vanishing fraction of the
vertices makes the two induced graphs agree; their common survivor is
\(Y\). It retains the constant-vector spaces almost isometrically and
has no short cycles. Proposition~\ref{prune:tower-properties} collects
these outputs. The finite alphabet must also remain small enough to
control the edges outside the later embedded copies. We begin with the
bit coordinates that make all these requirements compatible.

\subsection{Stream lengths and slots}

For \(h\geq0\), write \(h=2sq+t\), where \(q\geq0\) and
\(0\leq t<2s\), and set
\begin{equation}
\begin{split}
 l_A(h)&=r+n_0s+qs+\max(t-s,0),\\
 l_B(h)&=n_0s+qs+\min(t,s).
\end{split}
\label{tow:stream-lengths}
\end{equation}
We also write \(l_A(i)=l_A(|i|)\) and \(l_B(i)=l_B(|i|)\).
Thus \(l_A(h)+l_B(h)=r+2n_0s+h\).  Increasing \(h\) by one increases
exactly one of the two lengths by one: first the \(B\)-length for
\(s\) consecutive steps, and then the \(A\)-length for \(s\) steps.
Let
\[
 \begin{split}
 A_i&=\mathbb F _2^{l_A(i)}\times
       \mathbb F _2^{r+1+l_B(i)},\\
 B_i&=\mathbb F _2^{l_B(i)}\times
       \mathbb F _2^{s-r+1+l_A(i)}.
 \end{split}
\]
The two coordinates of either vertex are denoted \((u,v)\).
Since \(2r-s=\log_2d\), their exact cardinalities are
\begin{equation}
 |A_i|=d\,2^{H_i},\qquad |B_i|=2^{H_i}.
 \label{tow:cardinalities}
\end{equation}
All four active stream lengths at layer \(i\) are at most \(H_i\).
For example, \(l_A(i)\leq r+n_0s+|i|<H_i\), and the other three
bounds follow from the displayed definitions.

Partition every \(A\)-stream into a first block of length \(r\), called
block \(0\), followed by blocks of length \(s\), called
\(1,2,\ldots\).  Every \(B\)-stream has blocks of length \(s\), also
called \(1,2,\ldots\).  Order the \(u\)-blocks from the two sides as
\begin{equation}
 A,0;\ B,1;\ A,1;\ B,2;\ A,2;\ \ldots.
 \label{tow:interleaving}
\end{equation}
For every \(|i-k|\leq1\), the active bits of \(u_A\) at \(A_i\) and
\(u_B\) at \(B_k\) form an initial segment of this order, including
possibly a proper initial part of its last block.  Indeed this is true
for equal absolute indices by \eqref{tow:stream-lengths}; adjacent
absolute indices differ by adding a single bit of the next available
block, and choosing one of the two side lengths from each of these
adjacent indices preserves the same property.  Consequently the final
cutoff has exactly one of the two descriptions
\begin{equation}
\begin{array}{ll}
 l_A(i)=r+ns,\quad l_B(k)=ns+t,&1\leq t\leq s,\\
 l_B(k)=(n+1)s,\quad l_A(i)=r+ns+t,&1\leq t\leq s.
\end{array}
\label{tow:two-cutoffs}
\end{equation}
These descriptions include cutoffs at block ends by taking \(t=s\).
The initial value of \(n_0\) ensures that the four blocks immediately
before every such final block are full blocks of length \(s\).

For a pair \((i,k)\) with \(|i-k|\leq1\), put
\begin{equation}
\begin{aligned}
 \alpha&=1+l_B(i)-l_B(k),&p&=r+\alpha,\\
 \beta&=1+l_A(k)-l_A(i),&p'&=s-r+\beta.
\end{aligned}
\label{tow:offsets}
\end{equation}
Here \(\alpha,\beta\in\{0,1,2\}\), and \(p>p'>0\) for large
\(\ell\).  There are \(d\) slots at each \(A_i\)-vertex for this
pair, indexed by \(a=(k-i,j)\), \(1\leq j\leq d\).  At a
\(B_k\)-vertex the slots are
\(c=(i-k,j,b)\), where \(b\in\mathbb F _2^{p-p'}\).
Thus the desired degree on that side is
\begin{equation}
 D_{ik}=d\,2^{p-p'}=D\,2^{|i|-|k|},
 \qquad D/2\leq D_{ik}\leq2D.
 \label{tow:degrees}
\end{equation}
The identity follows from
\(p-p'=2r-s+|i|-|k|\).  In particular,
\(d|A_i|=D_{ik}|B_k|\).  After all three adjacent pairs are included,
the degrees are \(3d\) on \(A_i\) and at most \(6D\) on \(B_i\).

Set \(K=32s\).  The \emph{family} of a pair is the tuple
\[
 (k-i,\operatorname{sgn}(i),\operatorname{sgn}(k),
                  |i|\bmod K,|k|\bmod K),
\]
where zero has its own sign.  There are at most \(27K^2\) families.
The quantities \(\alpha,\beta,p,p'\), the slot conventions, and the
phase of the final block are determined by this tuple.  For a family
that occurs more than once, both indices have the same nonzero sign.
Its consecutive occurrences increase both absolute indices by \(K\),
and hence both side lengths by \(K/2=16s\).  Their final blocks are
therefore separated by \(32\) positions in
\eqref{tow:interleaving}.  Families involving zero occur only once.

\subsection{The linear slot evaluations}

On each side and for each slot define a linear map \(\mathcal K\) on
padded streams and its evaluation
\begin{equation}
 Q=v+\mathcal K u.
 \label{tow:evaluation}
\end{equation}
On a diagonal block of length \(h\in\{r,s\}\), \(\mathcal K\)
is multiplication by a coefficient of \(\mathbb F _{2^h}\), after
choosing a linear identification with \(\mathbb F _2^h\).  Only
\(A\)-block \(0\) uses \(h=r\).  From each block it may also send a
linear image into the first \(\alpha\) bits of the following block
on side \(A\), or into the first \(\beta\) bits of the following
block on side \(B\).  There are no other off-diagonal terms.  These
images will be called spills.  The same diagonal coefficient and spill
matrix are used at every repeated block of length \(s\) for a given
family and slot.  The spill from \(A\)-block \(0\) may, and will, be
zero.

Once the preceding input block is known, remove its spills from two slot
evaluations on the current block; their difference then cancels the
\(v\)-coordinate.
If the diagonal coefficients differ, their difference is invertible and
recovers the current \(u\)-block, followed by the \(v\)-block.  This is
the local operation behind recovery from a nonbacktracking path word.

We enforce this distinctness with finite global coefficient pools.
On side \(A\), a slot name is \((\delta,j)\), with
\(\delta\in\{-1,0,1\}\) and \(1\leq j\leq d\), so there are
\(3d\) names.  On side \(B\), include in the slot name the offset,
serial number, bucket width and bucket value.  The width belongs to
\(\{\log_2d-1,\log_2d,\log_2d+1\}\), so there are at most
\((21/2)D\) names.  Choose separate injections of the two pools into
\(\mathbb F _{2^s}\), and an injection of the \(A\)-pool into
\(\mathbb F _{2^r}\).  They exist because
\(3d<2^r\) and \((21/2)D<2^s\) for large \(\ell\).
Use the resulting coefficients independently of height and family.
In particular, any two different incident slots at a vertex have
different diagonal coefficients, for every applicable block size.
Spills can additionally depend on the family.  Finite fields introduce
no group elements: they are used only for these maps of bit coordinates.
Their existence follows, for example, by taking the roots of
\(z^{2^h}-z\) in a splitting field over \(\mathbb F _2\); the roots
form a field and, since the derivative is \(-1\), there are exactly
\(2^h\) of them.

Expansion requires a second property: a uniform bound on the fraction of
slots inducing the same input functional from a nonzero output functional,
including on partial blocks.  The following choice of injections and spills
supplies the needed multiplicity bounds.

\begin{lemma}[Choice of coefficients and spills]
\label{tow:coefficient-choice}
The injections and spills can be chosen so that the following holds
for every family, separately on each side.  Let \(\mathcal I\) be
the set of slots for its one adjacent pair, let \(n_{\rm sl}=|\mathcal
I|\), and let
\[
 M_a:\mathbb F _2^s\longrightarrow\mathbb F _2^{s+\alpha}
 \quad\hbox{or}\quad
 M_c:\mathbb F _2^s\longrightarrow\mathbb F _2^{s+\beta}
\]
be diagonal multiplication followed by its spill.  For every nonzero
linear functional \(\zeta\) on the output space, every value of the
restricted input functional \(\zeta M_a\), or \(\zeta M_c\), occurs
on at most \(0.55n_{\rm sl}\) slots after restriction to any nonempty
initial segment of the \(s\) input bits.  On all \(s\) input bits,
the corresponding bound is \(10^{-3}n_{\rm sl}\).
\end{lemma}

\begin{proof}
Choose each injection uniformly among injections into its field, and,
for every family and slot, choose an independent uniform two-row
binary spill matrix, retaining just its first \(\alpha\) or \(\beta\)
rows.  Suppose first that the spill component of \(\zeta\) is
nonzero.  Conditional on all diagonal coefficients, the full input
functionals \(\zeta M_a\) are independent and uniform among the
\(2^s\) binary functionals: a nonzero linear combination of independent
uniform rows is uniform.  Their restrictions to the first input bit
are therefore independent fair bits.  If the spill component is zero,
the diagonal component is nonzero.  Multiplication identifies the
\(2^s\) field coefficients bijectively with the \(2^s\) input
functionals, so the full functionals are sampled without replacement,
and their first-bit values are sampled without replacement from an
equally divided population.

Both sampling rules obey the following form of Hoeffding's concentration
bound \cite{Hoeffding1963}; we include the short argument:
\[
 \mathbb P\{S-\mathbb ES\geq \varepsilon m\}
       \leq \exp(-c_\varepsilon m)
\]
for the number \(S\) of successes in \(m\) samples and each fixed
\(\varepsilon>0\); the same bound holds for the lower tail.
For independent zero-one variables it follows by exponential Markov
and the bound
\(\mathbb E\exp(t(X-\mathbb EX))\leq\exp(t^2/8)\).
For sampling without replacement, the exponential moment is no larger
than for independent samples with the same success proportion.  Indeed,
the mean product of any fixed number of distinct nonnegative population
entries is at most the corresponding power of their mean: replacing
two unequal entries by their average increases each elementary
symmetric sum, because their sum is unchanged and their product
increases.  Repeating these replacements proves the inequality.
Apply it to the entries \(e^{tX}\), for either sign of \(t\).

Apply these bounds to each of the two first-bit values with threshold
\(0.55m\).  Apply them also to each full functional value with
threshold \(10^{-3}m\); its expected proportion is \(2^{-s}\),
which is less than \(10^{-3}/2\) for large \(\ell\).  In the
without-replacement full-functional case the stronger deterministic
bound of one occurrence is available.  The number of tests is at most
\(2^{O(s)}(1+K)^{O(1)}\), whereas every test has
\(m\geq d\) and failure probability at most \(e^{-cd}\), for
an absolute \(c>0\).  Their sum tends to zero as \(\ell\to\infty\).
Choose an outcome for which all tests pass.  Restricting to more than
one input bit only splits first-bit multiplicity classes, proving all
the asserted prefix bounds simultaneously.
\end{proof}

The output of \(\mathcal K\) fits in the active \(v\)-length at
every pair.  In fact those lengths equal
\begin{equation}
 |v_A|=p+l_B(k),\qquad |v_B|=p'+l_A(i).
 \label{tow:output-lengths}
\end{equation}
In the first case of \eqref{tow:two-cutoffs}, the \(A\)-output ends
by \(r+ns+\alpha\) and the \(B\)-output by
\((n+1)s+\beta\).  In the second case they end by
\(r+(n+1)s+\alpha\) and \((n+1)s+\beta\), respectively.
These are bounded by \eqref{tow:output-lengths}, including the cases
\(\alpha=0\), \(\beta=0\), and \(t=s\).

\subsection{One padded permutation rule for each family}

A positive edge label is a triple consisting of a family, a serial
number \(j\), and a vector \(x\in\mathbb F _2^p\).  Write
\(x=b\Vert x'\), where \(b\) has length \(p-p'\) and \(x'\)
has length \(p'\).  The reverse orientation has the formal inverse
label.  A family and \(x\) hence determine both slots, once their
serial number is given.

We next define triangular transformations \(\mathcal F\) of
\(B\)-blocks and \(\mathcal J\) of \(A\)-blocks.  Their parameters
include the family, \(j\), and \(x\).  At each block the rule may
also depend on all input \(u\)-blocks, on both sides, strictly earlier
in \eqref{tow:interleaving}.  This is the complete list of keys.
No transformation depends on input bits in its own block or a later
block through its keys.

For each actual pair in the family, mark its last active block and
its active prefix length \(t\).  At this marked block choose, for
each possible tuple of preceding keys, a permutation of
\(\mathbb F _2^t\); apply it to the first \(t\) bits and leave
the other bits of the block unchanged.  In the universal rule, apply
this permutation only when the four immediately preceding full
\(s\)-blocks of input are not jointly zero.  If they are jointly
zero, use the identity.  At unmarked blocks also use the identity.
At a \(B\)-block the transformation contributes to \(\mathcal F\),
and at an \(A\)-block to \(\mathcal J\).

The marked positions of one family are at least \(32\) blocks apart,
so there is at most one rule at any position.  Any finite prefix uses
only finitely many tables.  On inputs padded at the actual cutoff,
every higher marked rule is inactive: the first higher one has four
previous blocks past the cutoff, and its gate is closed; the same is
then true at every later marked block.  All these preceding inputs
are zero.  Thus the rule preserves the active input lengths.  It is
triangular and invertible, since its gate and permutation key are
known before its current block is inverted.  These statements also
hold when one side's entire stream is fixed and the other is solved
successively.

The equations defining the endpoints of an edge are
\begin{equation}
 Q_A=x\Vert\mathcal F(u_B),\qquad
 Q_B=x'\Vert\mathcal J(u_A).
 \label{tow:edge-equations}
\end{equation}
They are equations of padded streams.  To solve them from a vertex
\(A_i\) and its slot \(a\), compute \(Q_A\), read its first
\(p\) bits as \(x\), and obtain \(c\) from its bucket.  Invert
\(\mathcal F\) successively to recover \(u_B\); the known
\(u_A\) supplies its other keys.  Then form \(Q_B\) by the second
equation and set \(v_B=Q_B+\mathcal K_cu_B\).  Conversely, from
\(B_k,c\), the first \(p'\) bits of \(Q_B\) give \(x'\), and
the bucket in \(c\) gives \(x\).  Successively invert
\(\mathcal J\), then compute \(v_A\) from the first equation.
The support bounds above ensure that the resulting streams have the
required lengths.  In inversion beyond the cutoff, the unmarked
blocks copy zeros until the next marked block; its preceding four
blocks are zero, so its gate is closed, and induction continues.
The two constructions are inverse: they solve the same triangular
equations with the same keys.

Let \(Y^\#\) be the union of the blocks so defined.  Every slot has
exactly one incident edge, giving the degrees
\eqref{tow:degrees}.  Multiedges are allowed and are regarded as
distinct slot incidences.  The labeling is \emph{folded}: at an
\(A\)-vertex two distinct slots differ in offset or serial number;
at a \(B\)-vertex they differ in these data or in bucket.  In every
case the labels differ.  All outgoing labels at \(A\) are positive,
and all outgoing labels at \(B\) are formal inverses, so no further
coincidences are possible.  With \(m_\ell\) denoting the positive
alphabet size, we have
\begin{equation}
 m_\ell\leq 27K^2 d\,2^{r+2}
      =O\bigl(d\,2^r(1+s)^2\bigr),\qquad
 \log_2(2m_\ell)=11\ell+O(\log\ell).
 \label{tow:alphabet}
\end{equation}

For the auxiliary graph \(Y^0\), in each actual pair use the same
rule except that its own last-block permutation is applied
unconditionally.  All rules at lower marked positions retain their
gates.  Higher rules are inactive on its padded inputs.  The slot
inversion argument still applies, so the auxiliary block has the same
degrees and label set, and is folded as well.

\begin{lemma}[The independent matching variables]
\label{tow:matching-clusters}
Fix an actual pair \((i,k)\) and every table at its lower marked
positions.  The current last-block tables can be sampled so that the
remaining random choices in its \(Y^0\)-block are independent uniform
perfect matchings between pairs of clusters, each cluster having
\(2^t\) slot incidences.  There is one matching for each tuple
consisting of the family, serial number, \(x\), and all preceding
\(u\)-blocks on both sides.  Every edge in each cluster pair is a
possible edge, and its matching is used only for that pair of clusters.
\end{lemma}

\begin{proof}
Fix \(a\) and \(x\).  Since adding \(\mathcal K_au_A\) to
\(v_A\) is invertible, the \(A_i\)-vertices with these data are
freely parametrized by \(u_A\) and the tail of \(Q_A\), of lengths
\(l_A(i)\) and \(l_B(k)\).  On the other side fix the resulting
\(c\) and \(x'\).  Its vertices are freely parametrized by
\(u_B\) and the tail of \(Q_B\), of lengths \(l_B(k)\) and
\(l_A(i)\).  Thus both parameter spaces have dimension
\(l_A(i)+l_B(k)\).

The previously fixed transformations bijectively identify all data
before the final block, successively in the interleaved order.
Fix these preceding \(u\)-values.  In the first case of
\eqref{tow:two-cutoffs}, the remaining \(B\)-incidences are
parametrized by the \(t\) final bits of \(u_B\), while the remaining
\(A\)-incidences are parametrized by the corresponding \(t\) bits
of the tail of \(Q_A\).  No extra data remain, since the fixed
preceding values and these final bits specify the whole parameter
space on each side.  The current table of \(\mathcal F\) is
precisely a bijection between these two sets of \(2^t\) values.
Its variation can change the computed \(v_B\), but that coordinate
is already determined by the other endpoint data; it is not another
independent parameter.  In the second case the roles are reversed:
the \(t\) final bits of \(u_A\) are matched to the corresponding
tail bits of \(Q_B\) by the current \(\mathcal J\)-table.
Sampling a uniform permutation for each distinct full key tuple gives
the stated independent uniform matchings.  Different tuples specify
disjoint clusters of slot incidences.  The bucket is already part of
\(x\), so it supplies neither an omitted key nor a further coupling.
\end{proof}

The coefficient maps are now fixed, but the permutation tables remain to
be chosen.  The next section proves that the current matchings can be
chosen to give a uniform incidence gap, whatever lower-cutoff tables have
already been fixed.  That result will allow a successive choice in each
family: later tables cannot alter an earlier auxiliary block, since their
gates are inactive on its padded streams.
